%%======================================================================
\section{What Can Humanoid Robots Do Today?}
\label{sec:systems}
%%======================================================================

The global humanoid robot field in 2024--2026 comprises two largely distinct
populations: commercial platforms designed for industrial labor and a set of
academic research systems whose provenance spans two decades of bipedal
locomotion research. The central empirical finding of this section is that no
confirmed clinical deployment of any full-size bipedal humanoid robot exists as
of March 2026~\cite{atar2025humanoids,cho2026humanoid,cunha2025humanoid}.

Table~\ref{tab:platforms} provides the full specification comparison for all
thirteen commercial and selected academic platforms discussed in this section.
The analysis below synthesizes capability patterns and clinical relevance across
the field rather than characterizing individual platforms.

\begin{table*}[t]
  \caption{Humanoid Platform Capability Comparison (March 2026). Ht is reported
  in centimeters. N/A = not publicly disclosed. \textsuperscript{$\dagger$}Research prototype only, not
  clinical deployment. \textsuperscript{*}Not bipedal under this paper's scope
  definition; included as deployment benchmark.}
  \label{tab:platforms}
  \resizebox{\linewidth}{!}{%
  \begin{tabular}{@{}llrlllll@{}}
    \toprule
    \textbf{Robot} & \textbf{Type} & \textbf{Ht (cm)} &
    \textbf{Total DOF} & \textbf{Payload} &
    \textbf{LLM Integration} & \textbf{Open SW} &
    \textbf{Medical Deployment} \\
    \midrule
    Figure 02~\cite{figureai2024unveils} & Commercial & 168 & 28 & 20\,kg & Yes (in-house VLA)    & No      & None \\
    Tesla Optimus G2~\cite{tesla2023optimusgen2} & Commercial & 173 & 28+            & 20\,kg     & Yes (FSD-derived)     & No      & None \\
    BD Atlas (elec.)~\cite{bostondynamics2024atlas} & Commercial & $\sim$165 & 28      & 30\,kg & Yes (Gemini, announced) & No & None \\
    Agility Digit v4~\cite{agilityrobotics2023digit} & Commercial & 175 & N/A            & 16\,kg     & Announced             & Partial & None \\
    Unitree G1~\cite{unitree2024g1}       & Commercial & 132 & 23--43         & N/A        & Yes (VLA, open)       & Yes     & Research proto.~\cite{atar2025humanoids,cho2026humanoid}\textsuperscript{$\dagger$} \\
    Unitree H1~\cite{unitree2023h1}       & Commercial & 180 & 19             & N/A        & Yes (open)            & Yes     & None \\
    1X NEO Beta~\cite{onex2025neo}      & Commercial & 165 & 75             & 25\,kg     & Yes (proprietary)     & No      & None \\
    Apptronik Apollo~\cite{apptronik2023apollo} & Commercial & 173 & N/A            & 25\,kg     & Yes (Google AI)       & No      & None \\
    Sanctuary Ph. G8~\cite{sanctuaryai2023phoenix} & Commercial & 170 & N/A            & 25\,kg     & Yes (Carbon/Microsoft)& No      & None \\
    Fourier GR-2~\cite{fourier2024gr2}     & Commercial & 175 & 53             & N/A        & Announced             & Partial & None \\
    UBTECH Walker X~\cite{ubtech2024walkers}  & Commercial & 130--145 & 41        & 10\,kg    & Yes (multimodal LLM)  & No      & None \\
    AgiBot A2~\cite{agibot2025a2}        & Commercial & 175 & 49+            & N/A        & Yes (proprietary)     & No      & None \\
    NEURA 4NE-1~\cite{neura2025_4ne1}      & Commercial & $\sim$175 & N/A      & 15--20\,kg & Yes (cognitive AI)   & No      & None \\
    RHP Friends      & Academic   & 168 & 30             & N/A        & No                    & Partial & IREX 2023 nursing demo~\cite{benallegue2025rhp} \\
    NAO\textsuperscript{*} & Academic & 58 & 25         & N/A   & No                    & Partial & Human-Robot Interaction (HRI)/ASD trials (social only) \\
    Pepper\textsuperscript{*} & Academic & 120 & N/A   & N/A        & No                    & Partial & 2,000+ Japan care facilities \\
    \bottomrule
  \end{tabular}%
  }
\end{table*}

\subsection{Manipulation Maturity}
\label{sec:manipmaturity}

Dexterous manipulation is the primary functional gap between current humanoid
capability and clinical task requirements. Across the thirteen commercial platforms
surveyed, total articulation ranges from approximately 3~DOF per arm (Agility
Digit v4, designed for logistics bin-picking) to 75 total DOF (1X NEO Beta, with
22~DOF per hand). This spread reflects a fundamental architectural divergence:
logistics-oriented platforms optimize for payload and cycle time, whereas
platforms targeting service or research applications invest in hand complexity.

Payload capacity spans 10--30~kg across the commercial tier, well above the
requirements for most clinical manipulation tasks (stethoscope, ultrasound probe,
catheter, laparoscopic instrument). The binding constraint is not force output
but force \emph{control}: clinical contact requires sub-Newton precision at the
contact surface, a requirement that no platform has validated against patient-contact
standards such as ISO/TS~15066. The single exception is Sanctuary Phoenix Gen~8,
which specifies 5~mN tactile force sensitivity, the only commercial platform with
a disclosed haptic sensing specification relevant to clinical touch discrimination.

Documented medical procedure performance exists for one platform only: the
Unitree G1, in two independent studies achieving approximately 70\% aggregate
task success across seven clinical procedures under teleoperation~\cite{atar2025humanoids},
and demonstrating endoscope holding in a cadaveric surgical context~\cite{cho2026humanoid}.
Both results are preclinical prototypes. The cross-platform pattern is worth
noting: the most dexterous commercial platforms (NEO Beta, Sanctuary Phoenix)
have no clinical task evidence, while the only platform with clinical evidence
(Unitree G1) was chosen not for its dexterity but for its open SDK and research
accessibility. The bottleneck is therefore not raw manipulation capability but
manipulation that has been validated \emph{for clinical contact}.

\subsection{Locomotion and Clinical Environment Compatibility}
\label{sec:locomotion}

The 2022--2024 generation of commercial humanoids completed a decisive architectural
transition from hydraulic to electric actuation. All thirteen commercial platforms
surveyed are fully electric, with brushless DC or high-torque servo actuators.
This transition has three clinical implications: (1)~reduced acoustic output
enabling patient-proximate operation, (2)~elimination of fluid-system leak risk
in sterile environments, and (3)~simplified field maintenance without hydraulic
servicing. 1X NEO Beta specifies 22~dB operating noise (below the 30--35~dB
threshold typical of hospital wards at night), making it the only platform with
a disclosed acoustic specification relevant to patient care settings.

However, the clinical environment poses locomotion requirements that industrial
deployments have not exercised. Hospital flooring transitions between vinyl,
linoleum, ceramic tile, and threshold plates; patient room layouts require
operation in spaces as narrow as 90~cm; and patient contact during transfer
assistance imposes dynamic disturbance loads not present in logistics contexts.
No commercial platform has published locomotion performance data in hospital
environments, and no platform has been evaluated against IEC~62133 or IEC~60601
electrical safety standards in clinical settings.

Platform scale is a further structural constraint. The scope of this review sets
a 150~cm lower bound for human-adult-care tasks; of the 13 commercial platforms,
UBTECH Walker X (130--145~cm) and Unitree G1 (132~cm) fall at or below this
threshold. The Unitree G1 is retained as the only platform with clinical evidence
despite this limitation. Academic platforms designed for healthcare, including
RHP Friends (168~cm, 54~kg, 30~DOF), demonstrate that human-scale form factors
can be achieved with explicit clinical design intent~\cite{benallegue2025rhp}.

\subsection{LLM/VLA Integration and Task Architecture}
\label{sec:llmvla}

Twelve of the thirteen commercial platforms surveyed include announced or
deployed LLM/VLA integration~\cite{cao2025humanoid}. Architectures span three
patterns: (1)~in-house VLA foundation models trained on proprietary demonstration
data (Figure~02, Tesla Optimus); (2)~integration of open third-party models such
as Google Gemini Robotics (Boston Dynamics Atlas), open VLA frameworks
(Unitree G1 with LeRobot and OpenVLA support), or Microsoft-partnered AI
(Sanctuary Phoenix via Carbon); and (3)~multimodal LLM integration for
instruction following without explicit VLA policy training (UBTECH Walker X).

This architectural diversity obscures a uniform empirical gap: no VLA or LLM
system deployed on any humanoid platform has been trained on, evaluated against,
or validated for clinical task performance. The largest open VLA training corpus
(Open X-Embodiment, approximately 970,000 demonstrations) consists entirely of
household tabletop manipulation and pick-and-place tasks. Medical procedures
require physically and procedurally distinct demonstration data (instrument
grasps under sterile technique, force profiles during tissue contact, workflow
sequencing under infection control constraints), none of which appear in any
published training corpus.

The implication for clinical deployment is concrete: LLM/VLA integration as
currently implemented provides natural-language task specification and general
manipulation priors, but does not constitute a path to autonomous clinical task
execution. The Unitree G1's teleoperation results~\cite{atar2025humanoids,cho2026humanoid}
are informative because they bypass VLA autonomy entirely: the
cognitive and dexterous work is performed by a remote human operator, with the
robot serving as a physical proxy. Universal and dexterous human-to-humanoid
whole-body teleoperation has been demonstrated to enable complex bimanual
tasks~\cite{he2024omnikh}, consistent with the teleoperation-first architecture
these clinical studies employ. This teleoperation-first architecture may
represent the viable near-term clinical deployment model for all platforms,
regardless of their onboard AI sophistication.

\subsection{Platform Accessibility for Clinical Research}
\label{sec:accessibility}

Research access to humanoid platforms is heavily asymmetric. Unitree Robotics
is the only major commercial manufacturer to release and maintain a fully open
hardware SDK, which directly explains why the Unitree G1 is the only commercial
platform with peer-reviewed clinical task evidence. All other commercial platforms
operate under proprietary software stacks with no published API access for
external researchers. Academic platforms, including RHP Friends (30~DOF, purpose-built
for nursing) and iCub3 (54~DOF, full-body tactile skin, 40+ units in research
labs), provide broader research access but lack the scale and industrial backing
needed for clinical deployment pathways.

The open-source software ecosystem relevant to humanoid clinical research is
mature in general manipulation infrastructure but entirely absent in clinical-specific
components. A systematic survey of prominent repositories yields four layers:

\textbf{Foundation models and VLA frameworks.} LeRobot~\cite{cadene2024lerobot}
(HuggingFace; 22,179~stars; Apache~2.0) provides a hardware-agnostic imitation
learning and VLA pipeline with direct Unitree G1 support. Physical Intelligence's
openpi provides the $\pi_0$~\cite{black2024pi0} VLA model pre-trained on over
10,000~hours of real robot demonstrations, with best-reported results on ALOHA
bimanual benchmarks~\cite{zhao2023act}. OpenVLA~\cite{kim2024openvla} (7B parameters,
970,000 real-world demonstrations from Open X-Embodiment~\cite{openxembodiment2023};
MIT license) and Diffusion Policy~\cite{chi2023diffusion} complete the general-purpose
manipulation toolkit. All are general-purpose; none include healthcare-specific
training data, safety validation, or clinical task benchmarks.

\textbf{Simulation environments.} Genesis (28,254~stars; Apache~2.0) provides
unified differentiable physics simulation; Isaac Lab (6,544~stars; BSD-3-Clause)
offers GPU-accelerated reinforcement learning with native Unitree G1/H1 support.
Neither provides medical environment models, clinical workflow scenarios, or
validated patient-contact physics.

\textbf{Unitree hardware ecosystem.} The Unitree GitHub organization maintains
nine publicly accessible repositories including unitree\_rl\_gym (3,022~stars),
unitree\_sdk2 (932~stars), and xr\_teleoperate (1,309~stars) for full-body
XR teleoperation, the software component used in Atar et al.~\cite{atar2025humanoids}.
Low-cost leader-follower teleoperation hardware such as GELLO~\cite{wu2023gello}
further lowers the barrier to clinical procedure research.

\textbf{The clinical infrastructure gap.} Four categories of infrastructure
required for clinical deployment are absent from the open-source ecosystem:
(1)~a validated safety layer enforcing patient-contact force limits at runtime;
(2)~clinical workflow integration connecting humanoid actions to hospital EHR
and nursing systems; (3)~an open clinical demonstration dataset (every major
VLA training corpus consists of household and tabletop tasks); and (4)~a validated
medical simulation environment with clinical instrument physics and patient
anatomy proxies. The transition from general-purpose manipulation to clinically
safe, task-verified medical assistance is primarily a software and validation
challenge that the open-source community has not yet begun to address.

Table~\ref{tab:opensource} summarizes the most prominent open-source repositories
in this ecosystem.


\begin{table*}[t]
  \caption{Open-Source Software Ecosystem for Humanoid and Embodied-AI Robotics
  (March 2026). Stars are approximate GitHub counts at time of writing. Medical
  Relevance: \textbf{H}~=~high (direct hardware/algorithm support for clinical
  candidate platforms), \textbf{M}~=~medium (transferable general capability),
  \textbf{L}~=~low (infrastructure/simulation only). No repository provides
  clinical-specific safety validation or medical workflow integration.}
  \label{tab:opensource}
  \resizebox{\linewidth}{!}{%
  \begin{tabular}{@{}llrllp{4.5cm}l@{}}
    \toprule
    \textbf{Repository} & \textbf{Org / Author} & \textbf{Stars} &
    \textbf{Purpose} & \textbf{License} & \textbf{Notes} & \textbf{Med.\ Rel.} \\
    \midrule
    lerobot            & HuggingFace        & $\sim$22k  & VLA/imitation learning pipeline  & Apache 2.0  & Supports Unitree G1; no clinical data & H \\
    openpi             & Physical Intelligence & $\sim$10k & $\pi_0$ VLA foundation model      & Apache 2.0  & 10k+ hrs robot demos; bimanual ALOHA   & H \\
    openvla            & Stanford/Berkeley/TRI & $\sim$5.4k & 7B-param open VLA model          & MIT         & 970k demos; generalist manipulation    & H \\
    act                & tonyzhaozh         & $\sim$5k   & Action Chunking Transformers      & MIT         & Original ALOHA bimanual policy code    & H \\
    diffusion\_policy  & Columbia/MIT       & $\sim$4.5k & Diffusion-based robot policy      & MIT         & Imitation learning / VLA pipeline        & M \\
    Genesis            & Genesis-Embodied-AI & $\sim$28k & Differentiable physics simulation & Apache 2.0  & No medical environment models          & M \\
    IsaacLab           & isaac-sim (NVIDIA) & $\sim$6.5k & GPU-accelerated RL + sim          & BSD-3       & G1/H1 native support; no clinical env. & M \\
    unitree\_rl\_gym   & Unitree Robotics   & $\sim$3k   & RL locomotion for Unitree robots   & BSD-3       & G1/H1/B2 locomotion training           & H \\
    xr\_teleoperate    & Unitree Robotics   & $\sim$1.3k & Full-body XR teleoperation        & Apache 2.0  & Used in Atar et al.\ 2025 study        & H \\
    unitree\_sdk2      & Unitree Robotics   & $\sim$0.9k & Hardware SDK for G1/H1/B2         & Apache 2.0  & Low-level hardware interface           & H \\
    ManiSkill          & HKUST              & $\sim$1.8k & Embodied manipulation benchmark   & Apache 2.0  & No clinical task benchmarks            & M \\
    RoboManip          & isri-robot         & $\sim$0.3k & Bimanual manipulation dataset     & MIT         & Tabletop tasks only                    & M \\
    Open Duck Mini     & apirrone           & $\sim$2k   & Miniature bipedal robot (open HW) & MIT         & Educational; not clinical scale        & L \\
    \bottomrule
  \end{tabular}%
  }
\end{table*}

